\section{The Bailout Protocol}\label{sec:bailout}

The Bailout Protocol is the mandatory exception-escalation mechanism of the \bxthree{} Framework. It is not an error-handling routine selected at runtime, nor an optional escalation pathway available to a system that chooses to invoke it. It is an architectural compulsion: any unresolved exception state encountered by any node in a \bxthree{} system \emph{must} propagate upward through the accountability chain until it reaches a named human principal, and no machine actor along that chain may absorb, resolve, or suppress the exception. This section provides the complete formal specification: the deterministic state machine governing exception propagation, the precise conditions governing each transition, the bail-out trigger condition, the escalation algorithm, the bypass mechanism enforcing the no-machine-actor property, and the timeout handling that guarantees liveness under all failure conditions.

The protocol operates against the backdrop of the \bxthree{} three-layer model. The \purpose{} layer carries intent, judgment, and final authority; it defines the operating range $R(n)$ for each node $n$ and receives all escalated exceptions. The \bounds{} layer carries bounded reasoning and constrained execution within $R(n)$; it evaluates whether proposed actions fall within the node's provisioned authority and initiates bailout when they do not. The \fact{} layer carries deterministic enforcement and forensic auditability; it verifies proposed actions against the authorization ledger, enforces the state machine's transition relation, and maintains the append-only Forensic Ledger as an immutable record of every protocol event. Each node carries an immutable parent pointer $\alpha(n)$ established at provisioning and a reference to its human accountability anchor $h(n)$ — a named human principal — that is similarly immutable for the node's operational lifetime. The Bailout Protocol exploits this structure to guarantee that exception states propagate to $h(n)$ without passing through any intermediate machine actor as a terminus.

% --------------------------------------------------------------
\subsection{State Machine Definition}\label{sec:state-machine}
% --------------------------------------------------------------

The Bailout Protocol is formally modeled as a deterministic finite state machine embedded in the \fact{} layer of every \bxthree{} node:

\[
\mathcal{B} \;=\; (Q,\; q_0,\; \Sigma,\; \rightarrow,\; Q_F)
\]

\begin{itemize}\itemsep=0pt
\item[$Q$] $\{\;\text{IDLE},\; \text{BOUNDED},\; \text{BAIL\_PENDING},\; \text{ESCALATING},\; \text{HUMAN\_ACKNOWLEDGED},\; \text{RESOLVED}\; \}$ is the finite set of protocol states.
\item[$q_0$] $\text{IDLE}$ is the designated initial state, entered at node startup and after every successful directive resolution.
\item[$\Sigma$] The input alphabet comprising trigger events $e_{\text{trigger}}$, authorization confirmations $e_{\text{auth}}$, timeout signals $e_{\text{timeout}}$, human acknowledgments $e_{\text{ack}}$, binding resolution directives $e_{\text{resolve}}$, and rejection directives $e_{\text{reject}}$.
\item[$\rightarrow$] $\; \subseteq Q \times \Sigma \times Q$ is the deterministic transition relation (Table~\ref{tab:bailout-transition}).
\item[$Q_F$] $\{\text{RESOLVED}\}$ is the set of terminal accepting states.
\end{itemize}

A node's state $q \in Q$ is stored in the node's \fact{} layer worksheet and is observable by its parent node and by the Bailout Monitor. State updates are atomic: no intermediate or partially-executing states are observable externally. The machine is \emph{deterministic}: for every $(q, \sigma) \in Q \times \Sigma$, at most one $q' \in Q$ satisfies $(q, \sigma, q') \in \rightarrow$. Determinism is enforced by the Bailout Monitor's priority function $\pi$, which totally orders simultaneous trigger conditions and selects the highest-priority enabled transition before committing any state change. If no transition is defined for a $(q, \sigma)$ pair, the machine stalls and the Bailout Monitor records a protocol violation in the Forensic Ledger.

\begin{table}[htbp]
\caption{Bailout Protocol state transition relation $\rightarrow \subseteq Q \times \Sigma \times Q$. Rows are current state; columns are input events. Entries denote the next state. Blank cells indicate that no transition is defined — the machine stalls and the Bailout Monitor records a protocol violation.}
\label{tab:bailout-transition}
\begin{tabular}{@{}lcccccc@{}}
\toprule
 & \multicolumn{6}{c}{Input event} \\
\cmidrule(l){2-7}
Current state & $e_{\text{trigger}}$ & $e_{\text{auth}}$ & $e_{\text{timeout}}$ & $e_{\text{ack}}$ & $e_{\text{resolve}}$ & $e_{\text{reject}}$ \\
\midrule
IDLE              & BOUNDED            & — & — & — & — & — \\
BOUNDED           & — & HUMAN\_ACKNOWLEDGED & BAIL\_PENDING & — & — & — \\
BAIL\_PENDING     & — & — & ESCALATING & — & — & — \\
ESCALATING        & — & HUMAN\_ACKNOWLEDGED & ESCALATING & IDLE & RESOLVED & RESOLVED \\
HUMAN\_ACKNOWLEDGED & — & — & — & IDLE & RESOLVED & IDLE \\
RESOLVED          & — & — & — & — & — & — \\
\bottomrule
\end{tabular}
\end{table}

\bigskip
\noindent\textbf{State definitions.}

\begin{itemize}\itemsep=0pt
\item[\textbf{IDLE.}] The node is operating within its provisioned operating range $R(n)$. No exception condition is active. The \fact{} layer dispatches normal action requests through the standard execution path. The node may transition to BOUNDED upon receipt of a trigger event satisfying the bail-out trigger condition (TC), defined in \S~\ref{sec:bailout-trigger}.

\item[\textbf{BOUNDED.}] An exception $e$ has been detected by the \bounds{} layer and is being held at the \bounds{}--\fact{} layer boundary. A time-bounded resolution attempt is in progress. All resolution paths within current authority remain open. If the Bounds Engine produces a resolution $r$ that the \fact{} layer approves before $T_{\text{bounds}}$ expires, the node returns to IDLE. If all resolution paths are exhausted or $T_{\text{bounds}}$ expires, the state advances to BAIL\_PENDING.

\item[\textbf{BAIL\_PENDING.}] The node has determined — either through an explicit \texttt{bailout\_signal} or through $T_{\text{bounds}}$ expiration without a \fact-layer-approved resolution — that the exception exceeds its current authority. The Bailout Protocol is formally activated. This is the last architectural point at which a machine actor could theoretically suppress escalation. It cannot: entry to BAIL\_PENDING triggers an immediate atomic transition to ESCALATING via $e_{\text{timeout}}$ (T4), with no dwell time and no machine-actor intervention permitted.

\item[\textbf{ESCALATING.}] The exception is propagating upward toward $h(n)$. Each intermediate node receives the exception, appends its diagnostic context to the propagation chain, and forwards the exception via \texttt{SendToParent} without attempting local resolution. The state remains ESCALATING until either (a) the human anchor acknowledges receipt ($e_{\text{ack}}$: T5), (b) the human anchor issues a binding directive ($e_{\text{resolve}}$ or $e_{\text{reject}}$: T7), or (c) all retry pathways timeout after $N_{\max}$ retries at $T_{\text{cap}}$.

\item[\textbf{HUMAN\_ACKNOWLEDGED.}] The human anchor has received the exception, observed the diagnostic context, and acknowledged it. The node awaits a directive. No machine actor may issue, simulate, or substitute a directive in this state. The node remains until the human principal issues \texttt{ACK} (return to IDLE) or a binding directive \texttt{RESOLVE} / \texttt{REJECT} (enter RESOLVED).

\item[\textbf{RESOLVED.}] The human principal has issued a directive and the protocol has recorded it in the authorization ledger. The directive is one of: \texttt{ACK} (continue with current parameters), \texttt{RESOLVE} (execute a specific binding resolution), or \texttt{REJECT} (disallow the proposed action and enter a quiescent error state). This state is terminal for the current protocol run.
\end{itemize}

\bigskip
\noindent\textbf{State invariants.} The following invariants hold for all reachable states of $\mathcal{B}$:

\begin{align}
\text{INV}_1 &: \forall n \in \text{nodes},\; \text{state}(n) \in Q \setminus \{\text{ESCALATING}\} \iff \neg\text{active\_bailout}(n). \tag{INV-1}
\end{align}

That is, a node occupies a non-escalating state \emph{iff} it has no active bailout propagation in flight. A node may carry a bailout history (a ledger entry recording a past propagation) while in IDLE, but no active propagation may be underway.

\begin{align}
\text{INV}_2 &: \text{state}(n) = \text{BAIL\_PENDING} \implies \text{active\_bailout}(n) \land \text{transition}(\text{BAIL\_PENDING}, e_{\text{timeout}}) = \text{ESCALATING}. \tag{INV-2}
\end{align}

That is, any node in BAIL\_PENDING has an active bailout and the only permissible next transition is to ESCALATING, enforced by the Bailout Monitor without machine-actor discretion.

\begin{align}
\text{INV}_3 &: \text{state}(n) = \text{HUMAN\_ACKNOWLEDGED} \implies \text{directive\_source}(n) = h(n). \tag{INV-3}
\end{align}

No machine actor may issue a directive in HUMAN\_ACKNOWLEDGED; the only valid directive source is the designated human anchor $h(n)$.

% --------------------------------------------------------------
\subsection{Transition Conditions}\label{sec:transitions}
% --------------------------------------------------------------

\bigskip
\noindent\textbf{T1: IDLE $\rightarrow$ BOUNDED (trigger fires).} Transition T1 fires when the bail-out trigger condition TC holds for node $n$:

\begin{equation}
\text{T1 fires} \iff \text{TRIGGER}(n, x) \tag{T1}
\end{equation}

where $\text{TRIGGER}(n, x)$ is defined in \S~\ref{sec:bailout-trigger}.

\bigskip
\noindent\textbf{T2: BOUNDED $\rightarrow$ IDLE (local resolution approved).} T2 fires when the Bounds Engine produces a resolution $r$ that the \fact{} layer approves within $T_{\text{bounds}}$:

\begin{equation}
\text{T2 fires} \iff \text{state} = \text{BOUNDED} \land \text{FactLayer.approve}(r) \land \text{clock} < T_{\text{bounds}}. \tag{T2}
\end{equation}

\bigskip
\noindent\textbf{T3: BOUNDED $\rightarrow$ BAIL\_PENDING (bounds timeout).} T3 fires when $T_{\text{bounds}}$ expires without a \fact-layer-approved resolution:

\begin{equation}
\text{T3 fires} \iff \text{state} = \text{BOUNDED} \land \text{clock} \geq T_{\text{bounds}} \land \neg\text{FactLayer.approved\_resolution}. \tag{T3}
\end{equation}

\bigskip
\noindent\textbf{T4: BAIL\_PENDING $\rightarrow$ ESCALATING (compulsory escalation).} T4 fires immediately and atomically upon entry to BAIL\_PENDING, with no machine-actor discretion permitted:

\begin{equation}
\text{T4 fires} \iff \text{state} = \text{BAIL\_PENDING}. \tag{T4}
\end{equation}

\bigskip
\noindent\textbf{T5: ESCALATING $\rightarrow$ HUMAN\_ACKNOWLEDGED (human directive received).} T5 fires when the human anchor $h(n)$ issues a directive $d \in \{\texttt{ACK},\ \texttt{RESOLVE},\ \texttt{REJECT}\}$ and the \fact{} layer pre-commit hook verifies the directive originated from $h(n)$ (not from any machine actor):

\begin{equation}
\text{T5 fires} \iff d \in \{\texttt{ACK},\ \texttt{RESOLVE},\ \texttt{REJECT}\} \land \text{FactLayer.verifyOrigin}(d, h(n)). \tag{T5}
\end{equation}

If the directive is \texttt{ACK}, the node returns to IDLE (T8). If it is \texttt{RESOLVE} or \texttt{REJECT}, the protocol enters RESOLVED (T9).

\bigskip
\noindent\textbf{T6: ESCALATING $\rightarrow$ ESCALATING (retry on timeout).} T6 fires when the propagation timeout $T_{\text{prop}}$ expires without receiving an acknowledgment from the parent node:

\begin{equation}
\text{T6 fires} \iff \text{clock} \geq T_{\text{prop}} \land \text{attempts} < N_{\max}. \tag{T6}
\end{equation}

\bigskip
\noindent\textbf{T7: ESCALATING $\rightarrow$ RESOLVED (human unreachable).} T7 fires when all $N_{\max}$ retry attempts have been exhausted without receiving acknowledgment from the parent:

\begin{equation}
\text{T7 fires} \iff \text{attempts} \geq N_{\max} \land \neg\text{acknowledged}. \tag{T7}
\end{equation}

A \texttt{parent\_unreachable} event is recorded in the Forensic Ledger and propagated to $h(n)$ via the next available functional pathway.

\bigskip
\noindent\textbf{T8: HUMAN\_ACKNOWLEDGED $\rightarrow$ IDLE (acknowledgment without binding resolution).}

\begin{equation}
\text{T8 fires} \iff \text{state} = \text{HUMAN\_ACKNOWLEDGED} \land \text{directive} = \texttt{ACK}. \tag{T8}
\end{equation}

\bigskip
\noindent\textbf{T9: HUMAN\_ACKNOWLEDGED $\rightarrow$ RESOLVED (binding directive issued).}

\begin{equation}
\text{T9 fires} \iff \text{state} = \text{HUMAN\_ACKNOWLEDGED} \land \text{directive} \in \{\texttt{RESOLVE},\ \texttt{REJECT}\}. \tag{T9}
\end{equation}

% --------------------------------------------------------------
\subsection{Bail-Out Trigger Condition}\label{sec:bailout-trigger}
% --------------------------------------------------------------

The bail-out trigger is the condition that causes a node to enter BOUNDED state (T1). Rather than requiring system designers to enumerate specific error codes in advance — a practice that provides coverage only for anticipated failure modes and leaves unanticipated ones invisible — the Bailout Protocol defines the trigger in terms of \emph{observational boundary violation}: a node escalates precisely when its observable input space contains a condition outside its provisioned operating range $R(n)$ and the node cannot produce a \fact-layer-approved resolution within $T_{\text{bounds}}$.

\bigskip
\noindent\textbf{Primary trigger condition.} The trigger $\text{TRIGGER}(n, x)$ is defined as a conjunction of three sub-conditions, all of which must hold simultaneously:

\begin{equation}
\text{TRIGGER}(n, x) \;\equiv\; \underbrace{x \notin R(n)}_{(i)} \;\land\; \underbrace{\neg\text{resolved}(n, x)}_{(ii)} \;\land\; \underbrace{\text{confidence}(n, x) < \tau}_{(iii)} \tag{TC}
\end{equation}

where:

\begin{itemize}\itemsep=0pt
\item[(i)] $x$ is an element of the node's observable input space that falls outside the provisioned operating range $R(n)$. This is the \emph{bound signal}: the node observes something it was not provisioned to handle autonomously.

\item[(ii)] $x$ has not been previously resolved by the node within the current operational session. This prevents re-trigger on known-exception inputs that have already received human adjudication.

\item[(iii)] $\text{confidence}(n, x) < \tau$ is the node's self-assessed confidence that it can produce a correct resolution for $x$ without human input, where $\tau$ is the provisioned confidence threshold. The confidence function $\text{confidence} : \text{node} \times \text{input} \rightarrow [0, 1]$ is defined at node provisioning and is \emph{not} adjustable at runtime by any machine actor. This prevents a machine actor from inflating its own confidence to suppress a warranted escalation.
\end{itemize}

\bigskip
\noindent\textbf{Trigger priority function.} When multiple trigger conditions fire simultaneously for node $n$, they are totally ordered by the priority function $\pi(n, x)$, which ensures deterministic transition selection:

\begin{equation}
\pi(n, x) \;=\; w_1 \cdot \text{severity}(x) \;+\; w_2 \cdot \bigl(1 - \text{confidence}(n, x)\bigr) \;+\; w_3 \cdot \text{distance}\bigl(n, h(n)\bigr) \tag{PF}
\end{equation}

where $\text{severity}(x) \in [0, 1]$ is the assigned severity of exception $x$ (higher values indicate greater potential impact), $\text{distance}(n, h(n))$ is the number of hops from $n$ to the human anchor in the parent chain, and $w_1, w_2, w_3$ are provisioned constants satisfying $w_1 > w_2 > w_3 \geq 0$. The highest-priority trigger fires first. Ordering is enforced by the Bailout Monitor before any state transition is committed to the Forensic Ledger.

\bigskip
\noindent\textbf{Secondary trigger: explicit \texttt{bailout\_signal}.} In addition to TC, the Bounds Engine may emit an explicit \texttt{bailout\_signal} at any time, regardless of confidence scores. This handles the case where the Bounds Engine detects a condition — for example, a logical contradiction between two inputs within $R(n)$ — that it can recognize as unresolvable even though both inputs are individually within the operating range. The \texttt{bailout\_signal} bypasses TC entirely and immediately initiates T3:

\begin{equation}
\texttt{bailout\_signal}(x) \;\implies\; \text{T3 fires for node } n \text{ with exception } x. \tag{TC-secondary}
\end{equation}

\bigskip
\noindent\textbf{The no-discretion property.} The \bounds{} layer does not exercise discretion over the trigger condition. Specifically: (1) the evaluation of $\text{TRIGGER}(n, x)$ is a read-only query against the operating range definition $R(n)$, not a judgment call made by a machine actor; (2) the result of the query directly drives the state machine transition without an intervening decision step — there is no machine-actor gate between trigger evaluation and protocol activation; and (3) the \bounds{} layer may not delay, reclassify, or absorb an exception once $\text{TRIGGER}(n, x)$ evaluates to true.

% --------------------------------------------------------------
\subsection{Escalation Algorithm}\label{sec:escalation-algorithm}
% --------------------------------------------------------------

When the state machine enters ESCALATING (T4), the protocol executes the \texttt{Escalate} algorithm to propagate the exception to $h(n)$ along the parent chain. The algorithm is executed by the Bailout Monitor on the originating node; intermediate nodes execute only \texttt{ForwardException}, which appends diagnostic context and relays the exception upward.

\begin{algorithm}[htbp]
\caption{Bailout Protocol escalation algorithm.}
\label{alg:escalate}
\begin{algorithmic}[1]
\Procedure{Escalate}{$n$, trigger}
  \State $\text{ctx} \leftarrow$ \Call{BuildDiagnosticContext}{$n$, trigger}
  \State $p \leftarrow \alpha(n)$
  \State $\text{attempts} \leftarrow 0$
  \State $T_{\text{prop}} \leftarrow T_{\text{initial}}$
  \Repeat
    \State $\text{pathway} \leftarrow$ \Call{SelectPathway}{PHR, $p$}
    \State $\text{acked} \leftarrow$ \Call{SendToParent}{$p$, ctx, $T_{\text{prop}}$, pathway}
    \State \Call{UpdateLedger}{\text{propagation\_attempt}, $n$, $p$, $\text{pathway}$, $\text{clock}$}
    \If{$\text{acked}$}
      \State \Call{Transition}{\text{HUMAN\_ACKNOWLEDGED}}
      \State \Return \texttt{propagated}
    \EndIf
    \State $\text{attempts} \leftarrow \text{attempts} + 1$
    \State $T_{\text{prop}} \leftarrow \min(2 \cdot T_{\text{prop}},\ T_{\text{cap}})$
    \State \Call{UpdateLedger}{\text{retry}, \text{attempts}, \text{clock}$}
  \Until $\text{attempts} \geq N_{\max}$
  \State \Call{UpdateLedger}{\text{parent\_unreachable}, n, p, \text{clock}$}
  \State \Call{Transition}{\text{RESOLVED}}
  \State \Return \texttt{unresolved}
\EndProcedure
\\[6pt]
\Procedure{BuildDiagnosticContext}{$n$, trigger}
  \State $\text{ctx} \leftarrow \emptyset$
  \State $\text{ctx.trigger} \leftarrow trigger$
  \State $\text{ctx.node\_id} \leftarrow n$
  \State $\text{ctx.parent\_chain} \leftarrow ()$
  \State $\text{ctx.timestamp} \leftarrow \text{clock}()$
  \State $\text{ctx.propagation\_chain} \leftarrow []$
  \State \Return ctx
\EndProcedure
\\[6pt]
\Procedure{SelectPathway}{PHR, target}
  \State $\text{candidates} \leftarrow$ PHR.$\text{functional\_pathways\_to}$(target)
  \State \Return $\displaystyle\arg\min_{\text{pht}\in\text{candidates}} \text{pht.estimated\_latency}$
\EndProcedure
\end{algorithmic}
\end{algorithm}

\bigskip
\noindent\textbf{Propagation chain integrity.} The diagnostic context $\text{ctx}$ is immutable after construction by \texttt{BuildDiagnosticContext}. No node along the propagation path may modify, truncate, redact, or suppress any field of $\text{ctx}$. Each node appends a record to $\text{ctx.propagation\_chain}$ containing its node identifier, a timestamp, a state snapshot, and any locally relevant diagnostic metadata. This append-only integrity property is enforced by the \fact{} layer: the \texttt{UpdateLedger} call is atomic, and any attempt to modify an existing ledger entry is rejected and logged as a protocol violation.

\bigskip
\noindent\textbf{Correctness properties of \texttt{Escalate}.}

\begin{enumerate}
\item[\emph{Progress.}] If the parent node $p$ is reachable via at least one functional pathway, the algorithm eventually receives an ACK and returns \texttt{propagated}. The retry loop implements exponential backoff capped at $T_{\text{cap}}$, guaranteeing termination within a bounded number of attempts $N_{\max}$.

\item[\emph{Integrity.}] The diagnostic context is append-only and immutable after construction. No machine actor can alter, suppress, or selectively omit fields of the propagation chain. The \fact{} layer enforces this property at the ledger level: every \texttt{UpdateLedger} call is recorded, and the chain of records is cryptographically sealed.

\item[\emph{Liveness under pathway degradation.}] If $p$ is unreachable, the algorithm retries with exponential backoff. After $N_{\max}$ failures at $T_{\text{cap}}$, it fires a \texttt{parent\_unreachable} event, which propagates to $h(n)$ via the next available functional pathway tracked by the Pathway Health Registry. Pathways that fail $K$ consecutive health checks are marked \texttt{DEGRADED} and excluded from pathway selection until a health check succeeds.
\end{enumerate}

% --------------------------------------------------------------
\subsection{Bypass Mechanism}\label{sec:bypass-mechanism}
% --------------------------------------------------------------

The bypass mechanism enforces the \emph{no-machine-actor property}: the escalation from the Bailout Monitor to $h(n)$ must not traverse any machine actor that could absorb, redirect, or suppress the exception. This property is not a routing choice — it is a structural necessity.

\begin{definition}[Bypass Property]\label{def:bypass}
For every node $n \in \text{nodes}$, the set of permissible termination actors for an ESCALATING state is exactly $\{h(n)\}$:
\begin{equation}
\text{terminus}\bigl(\text{ESCALATING}(n)\bigr) \;=\; \{h(n)\}. \tag{BP}
\end{equation}
No machine actor $m \in \text{MachineActors} \setminus \{h(n)\}$ may appear in the termination set, may inspect the escalation payload, may modify the escalation, or may issue a directive that substitutes for one issued by $h(n)$.
\end{definition}

The bypass exists because any machine actor interposed between the Bailout Monitor and $h(n)$ introduces three mutually independent failure modes, each sufficient to break the upstream accountability guarantee:

\begin{enumerate}
\item[\textbf{Intercept.}] A machine actor could absorb the escalation, suppress the exception, and continue autonomous operation without ever reaching $h(n)$. This is precisely the failure mode the Bailout Protocol is designed to prevent.

\item[\textbf{Redirect.}] A machine actor could redirect the escalation to a different principal, creating an accountability gap in which the wrong human receives the exception and the correct human never learns of it.

\item[\textbf{Delay.}] A machine actor could buffer or queue the escalation, delaying $h(n)$'s receipt beyond $T_{\max}$ and violating the liveness guarantee.
\end{enumerate}

The bypass eliminates all three failure modes by architectural constraint: there is no machine actor in the communication path, so none of these failures are architecturally possible. This is analogous to a hardware interrupt line hardwired to a designated handler — just as a hardware interrupt cannot be intercepted or redirected by user-space code, the bailout escalation cannot be intercepted or redirected by any agent-level component.

The bypass is enforced by two cooperating mechanisms operating at different layers:

\begin{enumerate}
\item[\textbf{\fact{} layer gate.}] The \fact{} layer is the only mechanism capable of authorizing state transitions in the \bxthree{} Framework. The Bailout Monitor's pre-commit hook rejects any ledger entry for a resolution where the underlying trigger satisfies TC and $T_{\text{bounds}}$ has expired. This prevents the Bounds Engine from issuing a retroactive resolution after timeout to suppress warranted escalation. The \fact{} layer gate operates identically on all machine actors: there is no privileged machine actor with the ability to bypass this check.

\item[\textbf{Immutable parent-pointer and human-root architecture.}] The parent pointer $\alpha(n)$ and the human-root identifier $h(n)$ are stored in \fact-layer-managed memory and are read-only at the machine-actor level. A machine actor cannot redirect the parent chain, cannot insert itself as an intermediate resolver, and cannot modify the $humanRootId$ field that terminates the chain. Any attempted modification requires \purpose{} layer authorization with a corresponding ledger entry — which cannot be completed without human approval if the modification would affect the bailout terminus.
\end{enumerate}

Together, these mechanisms ensure that BP holds for all $n \in \text{nodes}$ under all execution conditions, including partial node failure, network partition, and deliberate attempt by a compromised machine actor to redirect the bailout pathway.

Conventional escalation hierarchies — including tiered oversight models that permit absorption at intermediate tiers — make human oversight a high-tier option rather than a compulsory terminus. The Bailout Protocol forbids this absorption by architectural constraint: $h(n)$ is not one option among many — it is the \emph{only} permissible terminus for every unresolved exception.

% --------------------------------------------------------------
\subsection{Timeout Handling}\label{sec:timeout-handling}
% --------------------------------------------------------------

Timeouts in the Bailout Protocol serve a liveness function: they prevent the protocol from permanently stalling when a parent node is unreachable, a communication pathway is unresponsive, or the human anchor is temporarily unavailable. Each timeout has a defined scope, a defined scope-exit action, and an escalation path that prevents deadlock.

\bigskip
\noindent\textbf{Timeout definitions.}

\begin{enumerate}
\item[\textbf{$T_{\text{bounds}}$ — Bounds Engine resolution timeout.}] The maximum time a node may remain in BOUNDED state attempting autonomous resolution. Default: 60 seconds, provisioned per node class according to the expected latency of the node's resolution procedure. If $T_{\text{bounds}}$ expires without a \fact-layer-approved resolution, the node transitions to BAIL\_PENDING $\rightarrow$ ESCALATING via T3 and T4. The \fact{} layer cannot extend $T_{\text{bounds}}$ unilaterally; any extension requires human authorization via a separate ledger entry recorded before $T_{\text{bounds}}$ expires.

\item[\textbf{$T_{\text{prop}}$ — Per-hop propagation timeout.}] The maximum time node $n$ waits for an acknowledgment from its parent after sending an exception via \texttt{SendToParent}. Default initial value: 5 seconds. If $T_{\text{prop}}$ expires without ACK, the algorithm retries with exponential backoff: $T_{\text{prop}} \leftarrow \min(2 \cdot T_{\text{prop}},\; T_{\text{cap}})$. This backoff continues until either an ACK is received or $N_{\max}$ retries have been exhausted at $T_{\text{cap}}$.

\item[\textbf{$T_{\text{cap}}$ — Maximum propagation timeout value.}] The ceiling value for the exponential backoff of $T_{\text{prop}}$. Default: 60 seconds. After $N_{\max}$ consecutive retries at $T_{\text{cap}}$ without ACK, the algorithm fires a \texttt{parent\_unreachable} timeout event, which propagates to $h(n)$ via the next available functional pathway tracked by the Pathway Health Registry.

\item[\textbf{$T_{\text{human}}$ — Human acknowledgment timeout.}] The maximum time the human anchor may remain in HUMAN\_ACKNOWLEDGED state without issuing a directive. Default: 300 seconds (5 minutes), configurable per deployment based on the expected human response time for the node's operational domain. If $T_{\text{human}}$ expires, the protocol issues a reminder notification to $h(n)$ and transitions the originating node to RESOLVED with an \texttt{unresolved} tag. The node becomes quiescent: no autonomous retry occurs, no further actions are attempted, and the \texttt{unresolved} tag in the Forensic Ledger signals that human re-engagement is required to clear the state. The protocol does not re-escalate automatically; automatic re-escalation would create a potential denial-of-service condition against the human anchor.
\end{enumerate}

\bigskip
\noindent\textbf{Pathway Health Registry.} The Pathway Health Registry (PHR) is a runtime data structure maintained by the Bailout Monitor that tracks the availability, latency, and error rate of each provisioned communication pathway to the human anchor. Before each \texttt{SelectPathway} call, the registry is queried and the pathway with the lowest estimated latency among currently functional pathways is selected. Pathways that fail $K$ consecutive health checks (default $K = 3$) are marked \texttt{DEGRADED} and excluded from pathway selection until a health check succeeds. A background health-ping thread updates the PHR at intervals of $T_{\text{health}}$ (default: 30 seconds), operating with the same scheduling priority as the Bailout Monitor to ensure that pathway status remains current even under heavy agent workload.

\bigskip
\noindent\textbf{Liveness guarantee.} Together, $T_{\text{bounds}}$, $T_{\text{prop}}$, $T_{\text{cap}}$, and $T_{\text{human}}$ define a finite upper bound on the wall-clock time from trigger firing (T1) to human acknowledgment under any combination of parent unavailability and pathway degradation. This bound is a deployment parameter — not a protocol constant — derived from the network topology, the number of hops to the human anchor, and the provisioned $T_{\max}$ parameter:

\begin{equation}
T_{\max} \;=\; T_{\text{bounds}} \;+\; N_{\max} \cdot T_{\text{cap}} \;+\; T_{\text{human}} \tag{T-max}
\end{equation}

A deployment that cannot guarantee human acknowledgment within the required $T_{\max}$ under all failure scenarios must provision additional or higher-bandwidth pathways, or must reduce $T_{\text{prop}}$ and $T_{\text{cap}}$ to bring the worst-case bound within the required limit. This analysis is a required part of deployment certification for any \bxthree{} system operating in a regulated or safety-critical domain.
